\section{The Resolver}
\label{sec:resolver}

The Zcash chain records Name Notes in canonical order, but it does not determine which transitions belong in the registry. A Resolver verifies each Name Note and applies accepted transitions to derive registry state.

Anyone can independently run a Resolver. Resolvers reading the same canonical Zcash chain under the current protocol rules must derive the same registry state.

\subsection{Resolver responsibilities}

A Resolver must:

\begin{enumerate}
  \item scan the canonical Zcash chain using the Mint's published full viewing key to identify and decrypt Ironwood outputs, obtaining their memos and note components;
  \item parse the decrypted outputs to identify Name Notes;
  \item verify each Name Note's commitment using the ZNS-specific derivation of $\mathsf{rcm}$ and $\psi$ defined in Section~\ref{sec:name-notes};
  \item verify Mint authorization and attestation as specified in Section~\ref{sec:the-mint};
  \item apply accepted transitions in canonical order to derive registry state; and
  \item provide clients with the current and historical binding state of each name, including the \field{expires\_at} value committed by the applicable Name Note.
\end{enumerate}

If a chain reorganization changes the canonical Zcash history, the Resolver MUST roll back any state derived from removed blocks and replay the replacement chain from the common ancestor.

\subsection{Verification}

Before state derivation, the Resolver verifies two properties of each Name Note:

\textbf{1. Commitment verification.} The Resolver parses the encoded transition and reconstructs the note commitment as specified in Section~\ref{sec:name-notes}. The Name Note passes commitment verification only if the reconstructed note commitment matches the \field{cmx} recorded on-chain.

\textbf{2. Mint authorization.} The Resolver uses the Mint's published viewing key to detect Name Note outputs. Detection alone does not establish that the Mint created an output: the viewing key is public, so any party can construct a well-formed Name Note output to the Name Note account. A Name Note is Mint-authorized only if its transaction spends a note that only the Mint spending key can spend: the current accepted Name Note of the name for an \action{update} or \action{release}, or a live anchor for a \action{claim}, as specified in Section~\ref{sec:resolver-state}. The Resolver also verifies the Mint's attestation as specified in Section~\ref{sec:the-mint}.

User authorization is not independently reconstructed by the Resolver. As described in Section~\ref{sec:authorization}, Request, Relay, and Respond are not part of the public Name Note history. Under the TEE assumptions, the Resolver relies on Mint attestation to establish that the approved authorization procedure was executed.

\subsection{State derivation}
\label{sec:resolver-state}

The Resolver processes transactions in canonical order: ascending block height, then transaction order within the block. A transaction carrying more than one Name Note produces no accepted transition.

For each name, the \field{prev\_rcm} field links a transition to the preceding Name Note within the current registration. A \action{claim}, including one following a \action{release}, uses the zero predecessor defined in Section~\ref{sec:name-notes}. An \action{update} or \action{release} is accepted only if its transaction spends the current accepted Name Note of that name and no other name's, and its \field{prev\_rcm} matches the \field{rcm} of that Name Note.

A \action{claim} has no predecessor Name Note to spend and instead spends an anchor: a zero-value note held by the Name Note account. The set of live anchors, the anchor pool, is established once during Mint setup with a fixed number of anchors. After the pool is established, other zero-value outputs to the Name Note account do not become anchors. A \action{claim} is accepted only if its transaction spends exactly one live anchor, spends no live Name Note, and creates exactly one zero-value output to the Name Note account, and the name has no live registration.

Every anchor a transaction spends leaves the pool, whether or not any transition is accepted. The zero-value output created by that transaction joins the pool as a successor anchor only if the transaction spent exactly one live anchor, carries exactly one Name Note, and creates exactly one zero-value output to the Name Note account. The Resolver derives the anchor pool from the canonical chain in the same order as registry state.

Because a Name Note can be spent only once, at most one transition can consume a given Name Note. Competing transitions that reference the same predecessor cannot all be accepted; the first to consume it on the canonical chain is the only one that can affect registry state.

A transaction that spends the current accepted Name Note of a name without creating a valid successor for that name ends the registration at that transaction. No \action{release} Name Note is created in this case.

Each accepted transition updates the registry state for that name. An \action{update} or \action{release} whose transaction does not spend the current accepted Name Note, and a \action{claim} whose transaction does not spend a live anchor, do not affect the registry.
